\documentclass[11pt]{article}

\usepackage{amsmath,amssymb,amstext,amsthm}
\usepackage{geometry}
\usepackage{fancyhdr}
\usepackage[pdftex]{graphicx}
\usepackage{xcolor}

\geometry{
  verbose,
  dvips,
  width=430pt, marginparsep=0pt, marginparwidth=0pt,
  top=90pt, headheight=12pt, headsep=20pt, footskip=30pt, bottom=80pt}

\setlength{\parskip}{\medskipamount}
\setlength{\parindent}{0pt}

\linespread{1.05}

\newtheorem{theorem}{Theorem}
\newtheorem{lemma}[theorem]{Lemma}
\newtheorem{proposition}[theorem]{Proposition}
\newtheorem{corollary}[theorem]{Corollary}
\newtheorem{conjecture}[theorem]{Conjecture}
\newtheorem{obs}{Observation}
\newtheorem{clm}{Claim}
\newtheorem{ques}{Question}
\newtheorem{problem}{Problem}

\newtheorem{ex}{Example}


\newcommand{\spc}{\hspace{0.08in}}
\newcommand{\vs}{\vspace*{.1in}}
\newcommand{\E}{\mathbb{E}}
\newcommand{\Prob}{\mathbb{P}}
\newcommand{\edit}[1]{\emph{\textcolor{blue}{#1}}}


\renewenvironment{proof}{{\noindent \textbf{\textit{Proof.}}}}
{\hfill $\Box$ \vspace*{0.1in}}
\newenvironment{proofc}{{\noindent \textit{Proof of Claim.}}}
{\hfill $\Box$}


\begin{document}

\begin{LARGE}\begin{center}\bf 14.3 Partial Derivatives\end{center}
\end{LARGE}

\vs\vs



\section{Warm Up}
\textbf{Question:} What does the derivative of a function of one variable represent?



\section{Motivation}
For a function of one variable (two dimensions on the plane $f(x)$ we have an unambiguous definition of a derivative. In 2D the derivative at a point $x$ is the slope of the tangent line to the curve at that point. The issue when we talk about surfaces in 3D is that there are infinite tangent lines to a surface at at point $(x,y)$ what do we do here? In this chapter we begin to answer this question.


\section{Definitions \& Theorems}

\begin{enumerate}
\item Let $f(x,y)$ be a function of two variables aka a surface. Since there are infinite tangent lines to the surface at $(x,y)$, we specifically set our direction to get one and only one line. Two directions which make sense to select is the $x$ direction and the $y$ direction. 

\item Since we want our tangent line to be in the direction of the $x$-axis, it must lie in a plane parallel to the $xz$ plane and cuts through our surface creating a curve. Any plane parallel to the $xz$ plane has $y$ constant. Thus, we call the derivative in the x direction the partial derivative with respect to $x$ and we denote this $f_x$ or $\partial f/ \partial x$. We think about $y$ in this case as a constant.

\item Similarly the partial derivative with respect to $y$ denoted $f_y$ or $\partial f/ \partial f$ is the tangent line to our surface in the direction of the $y$-axis. Here out line lies in a plane parallel to the $yz$ plane and we consider $x$ as a constant.

\item We can also write these derivative in terms of there limit definitions as follows.

\begin{equation*}
    f_x(x,y) = \lim_{h \to 0} \frac{f(x+h,y)-f(x,y)}{h}
\end{equation*}

\begin{equation*}
     f_y(x,y) = \lim_{h \to 0} \frac{f(x,y+h)-f(x,y)}{h}.
\end{equation*}

\item This definition can easily be extended beyond functions of two variable for instance if $w = f(x,y,z)$ is a function in $\mathbb{R}^4$ we have,

\begin{equation*}
   \partial w / \partial x =  f_x(x,y,z) = \lim_{h\to 0} \frac{f(x+h,y,z) - f(x,y,z)}{h}
\end{equation*}


\item We of course can take higher derivatives than just the first partial derivatives. The following figure 

\begin{figure}[h!]
    \centering
    \includegraphics[scale = .75]{Figures/partial der.png}
    \label{fig:enter-label}
\end{figure}

\item \begin{theorem}(Clairaut's Theorem)
    If $f$ is defined on a disk $D$ which contains $(a,b)$ and $f_x$ and $f_y$ are both continuous on $D$ then $f_{xy}(a,b) = f_{yx}(a,b)$.
\end{theorem}

\end{enumerate}




\section{Examples}

See partial-derivative worksheet


\newpage

\section{Scratch Work}

\end{document} 